%% kapa: exam=2 round=26 year=2015 period=1 subject=appraisal_practice subject_ko=감정평가실무 seq=3 number=3 points=20
%% source: Q-Net 감정평가사 기출문제 (cst00309) artlSeq=5149601 file=extracted/1008/5149601/2015년 제26회 감정평가사 2차시험 기출문제/1교시(감정평가실무).pdf
%% conversion: kordoc 4.24.51 -> tools/gen2.py (auto); review=proofread-gpt-sol; flags=none
\begin{longquestion}{【문제 3】}
감정평가사 김공정씨는 ○○ 택지개발사업지구로 지정고시된 지역의 보상에 대하여 중앙토지수용위원회로부터 이의재결 평가를 의뢰받았다. 관련법규를 참작하고 제시된 자료를 활용하여 다음의 물음에 답하시오. (20점) \par
\medskip\noindent 물음 1) 토지의 보상액을 감정평가하시오. (10점) \par
\medskip\noindent 물음 2-1) 관련법규에 의거 농업손실보상대상 여부를 검토하시오. (5점) \par
\medskip\noindent 2-2) 보상대상자별로 농업손실보상액을 산정하시오. (5점) \par
<자료 1> 사업의 개요 \par
{\small\centering
\resizebox{\ifdim\width>\linewidth\linewidth\else\width\fi}{!}{%
\begin{tabular}{|c|c|c|}
\hline
\makecell{사업의 종류 및 명칭} & \makecell{사업시행자} & \makecell{사업 위치 및 면적} \\
\hline
\makecell{○○ 택지개발사업} & \makecell{한국 △△ 공사} & \makecell{K도 P시 A동 및 Y동 일원 245,050 m\textsuperscript{2}} \\
\hline
\end{tabular}}\par}

1. 사업추진일정 \par
{\small\centering
\resizebox{\ifdim\width>\linewidth\linewidth\else\width\fi}{!}{%
\begin{tabular}{|c|c|}
\hline
\makecell{택지개발사업 주민등의 의견청취 공고일} & \makecell{2013. 06. 19} \\
\hline
\makecell{택지개발지구지정․고시일} & \makecell{2014. 01. 02} \\
\hline
\makecell{보상계획공고일} & \makecell{2015. 02. 05} \\
\hline
\makecell{재결일} & \makecell{2015. 08. 25} \\
\hline
\makecell{현장조사완료일} & \makecell{2015. 09. 19} \\
\hline
\makecell{이의재결일} & \makecell{2015. 10. 05} \\
\hline
\end{tabular}}\par}

<자료 2> 의뢰물건 내용 \par
1. 토지조서 \par
{\small\centering
\resizebox{\ifdim\width>\linewidth\linewidth\else\width\fi}{!}{%
\begin{tabular}{|c|c|c|c|c|c|}
\hline
\makecell{기호} & \makecell{소재지 지번} & \makecell{지목} & \makecell{공부면적(m\textsuperscript{2})} & \makecell{편입면적(m\textsuperscript{2})} & \makecell{비고(소유자)} \\
\hline
\makecell{1} & \makecell{P시 A동 10} & \makecell{임} & \makecell{1,200} & \makecell{1,200} & \makecell{이대한} \\
\hline
\end{tabular}}\par}

2. 물건조서 \par
{\small\centering
\resizebox{\ifdim\width>\linewidth\linewidth\else\width\fi}{!}{%
\begin{tabular}{|c|c|c|c|c|c|c|}
\hline
\makecell{기호} & \makecell{소재지 지번} & \makecell{물건의 유형} & \makecell{물건의 종류} & \makecell{구조 및 규격} & \makecell{수량} & \makecell{단위} \\
\hline
\makecell{1-1} & \makecell{P시 A동 10} & \makecell{농업손실 보상} & \makecell{당근} & \makecell{-} & \makecell{1} & \makecell{식} \\
\hline
\end{tabular}}\par}

<자료 3> 현장 조사내용 \par
1. <자료 2>의 기호(1) 토지는 토지소유자 이대한씨가 산지전용허가를 받지 아니하고 형질변경하여 경작해오다 건강악화로 김민국과 임대차계약서를 작성하고, 2013년 2월부터 김민국씨가 당근을 재배하고 있음 \par
2. 이대한씨와 김민국씨는 모두 해당지역에 거주하는「농지법」에서 정하는 농민으로, 농업보상에 대한 협의가 성립되지 아니한 상태임 \par
3. <자료 2>의 기호(1) 토지는 차량통행이 불가능한 노폭 약 2미터의 비포장도로에 접해 있고, 남서하향의 약 15도의 경사지에 위치한 부정형의 토지임. 한편, 용도지역은 계획관리지역임 \par
<자료 4> 표준지 공시지가 \par
{\small\centering
\resizebox{\ifdim\width>\linewidth\linewidth\else\width\fi}{!}{%
\begin{tabular}{|c|c|c|c|c|c|c|c|c|c|c|}
\hline
\multirow{2}{*}{\makecell{기호}} & \multirow{2}{*}{\makecell{소재지}} & \multirow{2}{*}{\makecell{면적\\(m\textsuperscript{2})}} & \multirow{2}{*}{\makecell{지목}} & \multirow{2}{*}{\makecell{이용\\상황}} & \multirow{2}{*}{\makecell{용도지역}} & \multirow{2}{*}{\makecell{도로교통}} & \multirow{2}{*}{\makecell{형상지세}} & \multicolumn{3}{|c|}{\makecell{공시지가(원/m\textsuperscript{2})}} \\
\cline{9-11}
 &  &  &  &  &  &  &  & \makecell{2013.\\01.01} & \makecell{2014.\\01.01} & \makecell{2015.\\01.01} \\
\hline
\makecell{가} & \makecell{P시 A동\\101} & \makecell{1,000} & \makecell{전} & \makecell{전} & \makecell{계획관리} & \makecell{세로\\(불)} & \makecell{세장형\\평지} & \makecell{68,000} & \makecell{74,000} & \makecell{77,000} \\
\hline
\makecell{나} & \makecell{P시 Y동\\15} & \makecell{1,000} & \makecell{임} & \makecell{자연림} & \makecell{계획관리} & \makecell{맹지} & \makecell{부정형\\완경사} & \makecell{21,000} & \makecell{23,000} & \makecell{24,000} \\
\hline
\makecell{다} & \makecell{P시 A동 산 11} & \makecell{12,000} & \makecell{임} & \makecell{자연림} & \makecell{생산\\관리} & \makecell{맹지} & \makecell{부정형\\급경사} & \makecell{15,000} & \makecell{17,000} & \makecell{18,000} \\
\hline
\end{tabular}}\par}

<자료 5> 시점수정 자료 \par
지가변동률(P시 계획관리지역, \%) \par
{\small\centering
\resizebox{\ifdim\width>\linewidth\linewidth\else\width\fi}{!}{%
\begin{tabular}{|c|c|}
\hline
\makecell{구분} & \makecell{변동률} \\
\hline
\makecell{2013. 01. 01 ～ 2013. 12. 31(누계)} & \makecell{4.213} \\
\hline
\makecell{2014. 01. 01 ～ 2014. 12. 31(누계)} & \makecell{2.765} \\
\hline
\makecell{2015. 01. 01 ～ 2015. 07. 31(누계)} & \makecell{1.175} \\
\hline
\makecell{2015. 07. 01 ～ 2015. 07. 31(당월)} & \makecell{0.167} \\
\hline
\end{tabular}}\par}

※ 2015년 8월 이후 지가변동률은 미고시 상태이며 2015년 7월 지가변동률과 동일하게 변동하는 것으로 추정함 \par
<자료 6> 표준지 공시지가 평균변동률(\%) \par
{\small\centering
\resizebox{\ifdim\width>\linewidth\linewidth\else\width\fi}{!}{%
\begin{tabular}{|c|c|c|c|}
\hline
\makecell{구분} & \makecell{2013년 ～ 2014년} & \makecell{2014년 ～ 2015년} & \makecell{2013년 ～ 2015년} \\
\hline
\makecell{K도} & \makecell{1.54} & \makecell{1.69} & \makecell{3.26} \\
\hline
\makecell{P시} & \makecell{3.51} & \makecell{4.23} & \makecell{7.89} \\
\hline
\end{tabular}}\par}

<자료 7> 지역요인 및 개별요인 비교자료 \par
1. K도 P시 A동 및 Y동은 인근지역임 \par
2. 도로접면 격차율 \par
{\small\centering
\resizebox{\ifdim\width>\linewidth\linewidth\else\width\fi}{!}{%
\begin{tabular}{|c|c|c|c|c|c|}
\hline
\makecell{광대세각} & \makecell{광대한면} & \makecell{소로한면} & \makecell{세로가} & \makecell{세로(불)} & \makecell{맹지} \\
\hline
\makecell{1.00} & \makecell{0.95} & \makecell{0.86} & \makecell{0.81} & \makecell{0.75} & \makecell{0.72} \\
\hline
\end{tabular}}\par}

3. 형상 격차율 \par
{\small\centering
\resizebox{\ifdim\width>\linewidth\linewidth\else\width\fi}{!}{%
\begin{tabular}{|c|c|c|c|}
\hline
\makecell{정방형} & \makecell{장방형} & \makecell{사다리형} & \makecell{부정형} \\
\hline
\makecell{1.00} & \makecell{0.99} & \makecell{0.98} & \makecell{0.95} \\
\hline
\end{tabular}}\par}

4. 지세 격차율 \par
{\small\centering
\resizebox{\ifdim\width>\linewidth\linewidth\else\width\fi}{!}{%
\begin{tabular}{|c|c|}
\hline
\makecell{평지} & \makecell{완경사} \\
\hline
\makecell{1.00} & \makecell{0.97} \\
\hline
\end{tabular}}\par}

<자료 8> 농업보상 자료 \par
1. 통계청 농가경제조사 통계자료 (도별 연간 농가평균 단위경작면적당 농작물 총수입) \par
(단위: 원, m\textsuperscript{2}) \par
{\small\centering
\resizebox{\ifdim\width>\linewidth\linewidth\else\width\fi}{!}{%
\begin{tabular}{|c|c|c|c|c|}
\hline
\makecell{행정\\구역} & \makecell{농작물\\총수입(원)} & \makecell{경지면적\\(m\textsuperscript{2})} & \makecell{농작물총수입/\\경지면적(원/m\textsuperscript{2})} & \makecell{2년분\\농업손실보상액(원/m\textsuperscript{2})} \\
\hline
\makecell{K도} & \makecell{18,855,086} & \makecell{11,086.12} & \makecell{1,701} & \makecell{3,402} \\
\hline
\end{tabular}}\par}

2. 실제소득인정기준에서 정하는 기관(농협)에서 발급받은 거래실적 자료 \par
{\small\centering
\resizebox{\ifdim\width>\linewidth\linewidth\else\width\fi}{!}{%
\begin{tabular}{|c|c|c|c|c|c|c|c|}
\hline
\makecell{출하주} & \makecell{출하처} & \makecell{품목} & \makecell{중량(kg)} & \makecell{평균판매단가\\(원/kg)} & \makecell{판매금액(원)} & \makecell{발급\\기관} & \makecell{비고} \\
\hline
\makecell{김민국} & \makecell{L마트\\외 4개소} & \makecell{당근} & \makecell{6,521} & \makecell{1,050} & \makecell{6,847,050} & \makecell{농협} & \makecell{연평균} \\
\hline
\end{tabular}}\par}

※ 김민국씨는 P시 A동 10번지 토지에서만 당근을 경작함 \par
3. 농축산물소득자료집 중 작목별 평균소득	(기준 연1기작/1,000 m\textsuperscript{2}) \par
{\small\centering
\resizebox{\ifdim\width>\linewidth\linewidth\else\width\fi}{!}{%
\begin{tabular}{|c|c|c|c|c|}
\hline
 & \makecell{수량(kg)} & \makecell{단가(원)} & \makecell{금액(원)} & \makecell{비고} \\
\hline
\makecell{조수입} & \makecell{4,184} & \makecell{832} & \makecell{3,481,088} & \makecell{-} \\
\hline
\makecell{생산비} & \makecell{-} & \makecell{-} & \makecell{1,595,346} & \makecell{종자비, 비료비 등} \\
\hline
\makecell{소득} & \makecell{-} & \makecell{-} & \makecell{1,885,742} & \makecell{소득률 54.2\%} \\
\hline
\end{tabular}}\par}

<자료 9> 기타자료 \par
1. <자료 4>의 기호(가)～(다)는 사업지구 내의 토지로서, 기호(가)～(다)의 평균변동률은 사업지구 내 표준지전체의 평균변동률과 동일함 \par
2. 그밖의 요인 보정: 대상토지의 인근지역 등의 정상적인 거래사례와 보상사례를 참작한 결과 농경지는 40 \% 상향보정, 임야는 80 \% 상향보정, 대지는 20 \%의 상향보정을 요함 \par
\end{longquestion}
